\originalchapter{简并量子气体与超流}
\label{ch:qfluid}
\source{8.333 L23--L26；理想气体先修见既有卷，重点为低能激发与响应。}
\originalsection{Fermi面与低温热激发}
单态占据$\bar n_k=[\e^{\beta(\epsilon_k-\mu)}+1]^{-1}$由$\Xi_k=1+z\e^{-\beta\epsilon_k}$直接得。零温费米面是占据突变的动量边界。低温只有其附近宽$T$的能量窗口可热激发，固定粒子数的Sommerfeld展开给$C/V=\pi^2D(E_F)\kb^2T/3$。一般晶格或各向异性色散，应以实际费米面$\int\dd S/[(2\pi)^d|\nabla\epsilon|]$计算态密度，不把自由气的球面公式当通用式。

相互作用Fermi液体的准粒子寿命率近$E^2+T^2$，其速度与热态密度可重整化；磁化率、压缩率还含Landau参数，详细构造见新凝聚态卷。统计力学中的“低温像自由准粒子”不意味着原相互作用消失，也不意味着所有观测都只需把$m$换成$m^*$。
\originalsection{Bose凝聚与超流的区别}
玻色占据$\bar n_k=[\e^{\beta(\epsilon_k-\mu)}-1]^{-1}$要求$\mu\le\epsilon_{\min}$。在三维抛物线气体，激发态能容纳的密度上限$n_{\mathrm{ex}}=\zeta(3/2)/\lambda_T^3$，超出的粒子进入最低态。均匀$d\le2$理想气有限温的低能积分发散，因此没有同样的有限温凝聚。有限势阱、相互作用与二维BKT需分别讨论。

凝聚是单体密度矩阵的一个本征值$O(N)$，超流是对相位扭曲、旋转或流动的刚度和耗散性质。理想Bose气可凝聚，却在$\epsilon(p)=p^2/2m$下$\min\epsilon/p=0$，无正Landau临界速度。弱排斥作用使Bogoliubov谱$E_p=\sqrt{(p^2/2m)(p^2/2m+2gn)}$低能线性，才给非零声速尺度。
\originalsection{Landau准则、声子与roton}
在速度$u$的流体中，创建一激发同时减少流体动量$p$。Galilean能量变化为$\Delta E=\epsilon(p)-u\cdot p$，若所有激发皆正则弱扰动不能自发激发它们，所以$u_c=\min_p\epsilon(p)/|p|$。$^4$He的声子线性段给$c$，roton谷$\epsilon(p)\simeq\Delta_r+(p-p_0)^2/(2m_r)$可给更小约$\Delta_r/p_0$。涡旋成核、边界与非线性过程可使实际临界速度更小，Landau准则不是任意实验的充分无耗散保证。

用密度激发试态$\ket q=\rho_q\ket0$，归一为$N S(q)$，双对易和规则给激发变分能
\begin{equation}
E_F(q)=\frac{\hbar^2q^2}{2mS(q)}.
\end{equation}
它是Feynman密度试态的能量上界而非所有激发的精确色散。静态结构因子峰表示密度排列的强相关，对应上界变低，解释roton与液体短程结构的联系。
\originalsection{正常密度与两流体}
激发在小相对速度$w$下分布$n_B(\epsilon_p-p\cdot w)$，其动量密度线性为$\rho_n w$，所以
\begin{equation}
\rho_n=\frac1d\int\frac{\dd^dp}{(2\pi\hbar)^d}p^2[-\partial_{\epsilon_p}n_B].
\end{equation}
三维声子$\epsilon=cp$给$\rho_n=2\pi^2(\kb T)^4/(45\hbar^3c^5)$，积分中$\int x^4\e^x/(\e^x-1)^2=4\pi^4/15$。roton附近按Gaussian径向能量积分贡献$\propto T^{-1/2}\e^{-\Delta_r/\kb T}$，数值前因子含$p_0,m_r$。超流密度$\rho_s=\rho-\rho_n$。

两流体总质量流$\rho_su_s+\rho_nu_n$，熵主要由正常部分携带。线性化密度与熵守恒可有第一声（主要密度波）与第二声（主要熵温度波），低热膨胀近似$u_2^2=(\rho_s/\rho_n)Ts^2/c_p$，此处$s,c_p$按单位质量。式中若改成每体积熵须同步换单位，否则量纲会错。两流体是低能局域平衡描述，不从“有BEC”一词自动推出。
